% !TEX root = ../main.tex
\chapter{AURORA IDE: arquitetura geral}
\label{cap:05-aurora-arquitetura}

A \tool{AURORA} é a IDE desktop da plataforma \tool{SAPHO}: uma aplicação
\href{https://www.electronjs.org/}{Electron} na qual o usuário projeta
soft-processors, edita código em \cmm, Verilog e SystemVerilog, compila pela
toolchain \tool{YANC} (\autoref{cap:07-yanc-pipeline}), simula, visualiza o RTL
no \tool{PRISM} e conversa com o subsistema Aurora Intelligence
(\autoref{cap:12-aurora-intelligence}). Este capítulo documenta a
\emph{arquitetura do processo Electron}: a separação entre processo principal
(\emph{main}) e processo de renderização (\emph{renderer}), as janelas criadas,
a topologia de IPC (\emph{inter-process communication}) e os módulos de
infraestrutura do \emph{main} (estado, ciclo de vida, caminhos, rastreio de
processos filhos, log, carregamento do bundle do renderer).

O escopo aqui é o processo \emph{main} e a fronteira IPC. A interface em si
(componentes \href{https://lit.dev/}{Lit}, editor \tool{Monaco}, terminal,
painéis) é tratada em \autoref{cap:06-aurora-renderer}; a toolchain, o pipeline
de compilação e o \tool{PRISM} têm capítulos próprios. Quando um assunto pertence
a outro capítulo, ele é apenas mencionado de passagem.

\begin{nota}
Toda afirmação deste capítulo foi verificada lendo o código-fonte em
\path{C:/Users/chrys/Documents/GitHub/aurora}. Os arquivos centrais são
\file{main.js} (ponto de entrada), \file{main/windows.js} (fábricas de janelas),
\file{main/lifecycle.js} (ciclo de vida), os módulos \file{main/ipc/*.js}
(handlers IPC) e os preloads em \path{js/app/preload*.js}.
\end{nota}

\section{Modelo de dois processos do Electron}
\label{sec:05-dois-processos}

O Electron herda do Chromium um modelo multiprocessos. A \tool{AURORA} usa duas
classes de processo:

\begin{description}[leftmargin=*]
  \item[Processo \emph{main} (Node.js)] É o processo privilegiado. Tem acesso a
    Node, ao sistema de arquivos, ao \lstinline|child_process| (spawn de
    compiladores/simuladores), à rede e às APIs nativas do Electron
    (\lstinline|app|, \lstinline|BrowserWindow|, \lstinline|dialog|,
    \lstinline|session|, \lstinline|safeStorage|). Existe exatamente um por
    instância da aplicação. Todo o código sob \file{main/} roda aqui.
  \item[Processo \emph{renderer} (Chromium)] É a página web que desenha a UI.
    Roda \emph{sandboxed}, sem acesso direto a Node nem ao filesystem. Cada
    \lstinline|BrowserWindow| tem o seu próprio renderer. Para alcançar qualquer
    capacidade do \emph{main}, o renderer chama uma API curada exposta por um
    \emph{preload} via \lstinline|contextBridge|.
\end{description}

\subsection{O padrão central: ``renderer orquestra, main executa''}
\label{sec:05-padrao-central}

A linha divisória de toda a \tool{AURORA} é: \emph{o renderer decide o que fazer;
o main faz o que pode causar efeito colateral}. O renderer monta a intenção
(abrir projeto, ler arquivo, compilar uma etapa, fazer commit, gerar um SVG) e a
envia por um canal IPC; o \emph{main} valida a entrada, executa a operação
privilegiada (filesystem, \lstinline|spawn|, rede, keychain) e devolve o
resultado. Nenhuma capacidade de Node, \lstinline|fs| ou \lstinline|spawn| vive
no renderer --- todas estão no \emph{main}, atrás de um canal IPC.

Esse padrão aparece de forma explícita no executor de compilação. O renderer
constrói um \lstinline|CommandSpec| estruturado (\lstinline|{binary, args}|) e o
manda por \lstinline|exec-spec|; o \emph{main} valida o binário contra uma
allowlist da toolchain, confere \emph{protected flags} e só então dispara
\lstinline|spawn(binary, args, { shell: false })|. Como descrito no cabeçalho de
\file{main/compile/executor.js}, com \lstinline|shell:false| cada argumento vai
para o filho exatamente como está --- não há parsing de shell, portanto não há
bug de \emph{quoting} nem vetor de injeção. É exatamente essa propriedade que
permite que o sistema de \emph{command override} da Aurora Intelligence exista
com segurança (\autoref{cap:12-aurora-intelligence}).

\section{O ponto de entrada: \file{main.js}}
\label{sec:05-mainjs}

O arquivo \file{main.js} é deliberadamente magro: toda a lógica real vive em
\file{main/}, e este arquivo apenas \emph{conecta os módulos} numa ordem de boot
determinística. Suas responsabilidades, em ordem de execução, são:

\begin{enumerate}[leftmargin=*]
  \item \textbf{Configurar o logger antes de tudo.} A primeira instrução após os
    \lstinline|require| é \lstinline|configureLogger()| (de \file{main/logger.js}),
    para que toda chamada de log subsequente já use a configuração definitiva.
  \item \textbf{Rede de segurança do processo main.} Inicia o
    \lstinline|crashReporter| em modo local (\lstinline|uploadToServer: false|,
    coletando minidumps) e registra \lstinline|process.on('uncaughtException')|
    e \lstinline|process.on('unhandledRejection')| para que nada que escape de um
    callback derrube o processo silenciosamente. Tudo \emph{best-effort}: a rede
    de segurança nunca pode ser o que bloqueia o boot.
  \item \textbf{Ajuste de GPU / compositor.} As \emph{command-line switches} têm
    de ser anexadas \emph{antes} de o app ficar pronto, pois o processo de GPU as
    lê no lançamento. Veja \autoref{sec:05-gpu}.
  \item \textbf{AppUserModelID (Windows).} Chama
    \lstinline|app.setAppUserModelId('com.nipscern.sapho')| o mais cedo possível
    --- antes de existir qualquer \lstinline|BrowserWindow| --- para que o Windows
    associe o processo a uma identidade estável de \emph{jumplist}/taskbar.
  \item \textbf{Registro de IPC e ciclo de vida.} Adquire o \emph{single-instance
    lock} via \lstinline|lifecycle.register()| e, se obtido, registra todos os
    módulos de IPC (\autoref{sec:05-ipc}).
  \item \textbf{\lstinline|app.whenReady()|.} Quando o Electron está pronto:
    cria o diretório \path{/tmp} do MSYS, roda a higiene de temporários
    (\file{main/temp_gc.js}), reconstrói a jumplist, instala a
    \emph{Content-Security-Policy} (\autoref{sec:05-csp}) e finalmente chama
    \lstinline|windows.createSplashScreen()| --- a splash agenda a criação da
    janela principal.
\end{enumerate}

\subsection{Switches de GPU}
\label{sec:05-gpu}

Antes de \lstinline|app.whenReady|, \file{main.js} anexa duas switches:

\begin{lstlisting}[language=Java]
app.commandLine.appendSwitch('enable-gpu-rasterization');
app.commandLine.appendSwitch('enable-zero-copy');
\end{lstlisting}

A combinação rasteriza os \emph{tiles} na GPU e os carrega sem cópia pela CPU ---
um ganho para painéis de muito \emph{scroll}/\emph{paint} (terminal, editor,
configuração de waveform). O código documenta explicitamente que \emph{não} se
usa \lstinline|disable-frame-rate-limit| (removeria o \emph{vsync} do Chromium e
poderia saturar o processo de GPU durante os \emph{loops} de
\lstinline|requestAnimationFrame| da splash). A aceleração de hardware é deixada
\textbf{ligada} --- \lstinline|app.disableHardwareAcceleration()| nunca é chamado
--- e o \lstinline|backgroundThrottling| fica no padrão.

\subsection{Content-Security-Policy}
\label{sec:05-csp}

Dentro de \lstinline|app.whenReady|, \file{main.js} instala uma CSP via
\lstinline|session.defaultSession.webRequest.onHeadersReceived|, entregue como
\emph{header} de resposta para que cubra tanto o carregamento empacotado
(\lstinline|file://|) quanto o servidor de desenvolvimento. O \emph{handler}
\emph{substitui} (não acrescenta) qualquer CSP que um servidor de dev tenha
definido, deixando a da \tool{AURORA} como autoritativa. As diretivas incluem,
entre outras, \lstinline|default-src 'self'|, \lstinline|object-src 'none'|,
\lstinline|frame-ancestors 'none'| e um \lstinline|connect-src| que, em produção,
só permite a origem própria mais a detecção local do \tool{Ollama}
(\lstinline|localhost:11434|) --- os provedores de IA em nuvem e o \tool{MCP}
rodam no \emph{main}, não no renderer.

\section{Janelas: \file{main/windows.js}}
\label{sec:05-janelas}

As fábricas de janela vivem em \file{main/windows.js}, com exceção da janela do
\tool{PRISM}, que mora em \file{main/ipc/prism.js} por carregar muita lógica de
compilação. Ao todo a \tool{AURORA} cria \textbf{cinco} tipos de
\lstinline|BrowserWindow|: principal, splash, atualização, \tool{PRISM} e o
\emph{Design Lab}. A tabela abaixo resume cada uma e suas opções de segurança ---
em todas elas \lstinline|contextIsolation| está ligado, \lstinline|sandbox| está
ligado e \lstinline|nodeIntegration| está desligado.

\begin{longtable}{@{}p{0.16\linewidth} p{0.30\linewidth} p{0.46\linewidth}@{}}
  \toprule
  \textbf{Janela} & \textbf{Preload dedicado} & \textbf{Características} \\
  \midrule
  \endhead
  Principal & \file{js/app/preload.js} &
    \lstinline|frame:false|, \lstinline|thickFrame:true|, barra de título
    customizada; 1280\texttimes820 (mín.\ 720\texttimes480); carrega
    \path{index.html}. Bloqueio de navegação: \lstinline|setWindowOpenHandler|
    nega \lstinline|window.open| e \lstinline|will-navigate| é cancelado.
    Múltiplas janelas principais podem coexistir no mesmo processo. \\
  Splash & \file{js/app/preload_splash.js} &
    720\texttimes480, \lstinline|frame:false|, \lstinline|transparent:true|,
    \lstinline|alwaysOnTop|, \lstinline|skipTaskbar|; carrega
    \path{html/splash.html}. Exibe progresso \emph{real} de boot. \\
  Atualização & \file{js/app/preload_update.js} &
    540\texttimes660, \lstinline|frame:false|, \lstinline|transparent:true|,
    \lstinline|alwaysOnTop|; janela única que percorre os estados
    \emph{available} / \emph{downloading} / \emph{downloaded}
    (\autoref{cap:13-updater-build}). Não fecha durante um download em curso. \\
  \tool{PRISM} & \file{js/app/preload_prism.js} &
    1400\texttimes900 (mín.\ 1000\texttimes700), \lstinline|frame:false|,
    \lstinline|thickFrame:true|; carrega \path{html/prism/prism.html}.
    Criada em \file{main/ipc/prism.js} (\autoref{cap:08-prism-rtl}). \\
  Design Lab & (sem preload) &
    1100\texttimes820, janela emoldurada (controles nativos do SO), singleton.
    Galeria interna de componentes, aberta pela paleta de comandos; é uma página
    estática que não fala com nada. \\
  \bottomrule
  \caption{As cinco janelas da \tool{AURORA} e suas opções principais.}
\end{longtable}

\subsection{Janela principal: \lstinline|createMainWindow|}
\label{sec:05-main-window}

A janela principal é criada sem moldura nativa (\lstinline|frame:false|) porque a
\tool{AURORA} desenha a própria barra de título; \lstinline|thickFrame:true|
preserva o \emph{Aero snap}, o redimensionamento de borda e as animações em todas
as versões do Windows. As \lstinline|webPreferences| materializam o modelo de
segurança: \lstinline|contextIsolation:true|, \lstinline|sandbox:true|,
\lstinline|nodeIntegration:false|, \lstinline|nodeIntegrationInSubFrames:false|,
\lstinline|webviewTag:false|, \lstinline|enableWebSQL:false|,
\lstinline|allowRunningInsecureContent:false| e
\lstinline|experimentalFeatures:false|. O único caminho do renderer para o
\emph{main} é o \lstinline|electronAPI| exposto pelo preload.

A função também:

\begin{itemize}[leftmargin=*]
  \item registra o protocolo \lstinline|sapho:| e a extensão \path{.spf} no
    Windows e reconstrói a \emph{jumplist} a cada criação de janela;
  \item notifica o renderer sobre maximizar/restaurar/fullscreen (evento
    \lstinline|window-state|) para que o ícone da barra de título acompanhe;
  \item ao detectar um \path{.spf} passado na linha de comando, envia
    \lstinline|open-spf-file| ao renderer após \lstinline|did-finish-load|;
  \item no \lstinline|close| da \emph{última} janela principal, dispara
    \lstinline|stopAllToolchain()| e fecha a janela auxiliar do \tool{PRISM},
    garantindo que nenhum compile/simulação/síntese sobreviva ao fechamento da
    IDE.
\end{itemize}

\subsection{Splash com progresso real}
\label{sec:05-splash}

A splash não é um \emph{timer} fixo: a barra avança conforme a janela principal
cruza marcos genuínos de boot. O fluxo, codificado em
\lstinline|createSplashScreen|, é:

\begin{lstlisting}[language=bash]
splash carrega (did-finish-load)         -> progresso 10%  "boot"
  cria janela principal OCULTA (deferShow)
                                         -> progresso 24%  "window"
                                         -> progresso 36%  "loading"
  webContents da principal:
    dom-ready                            -> progresso 56%  "dom"
    did-finish-load                      -> progresso 80%  "resources"
    ready-to-show                        -> progresso 90%  "editor"
  renderer sinaliza app:renderer-ready   -> progresso 96%  -> handoff 100%
  splash responde splash:filled          -> espera 2 s     -> reveal
\end{lstlisting}

O \lstinline|reveal| maximiza e mostra a janela principal e fecha a splash de
imediato (sem \emph{fade}, para não haver \emph{flash} de desktop vazio entre
elas). Um teto absoluto de \textbf{15\,s} força a transferência caso algum marco
nunca dispare. O coordenador usa \lstinline|ipcMain.once| para os sinais
\lstinline|splash:filled| e \lstinline|app:renderer-ready|, de modo que uma
segunda janela criada depois não re-dispare o \emph{handoff}.

\subsection{Controles de janela roteados pelo emissor}
\label{sec:05-window-controls}

Como vários janelas principais podem coexistir no mesmo processo (ver
\autoref{sec:05-lifecycle}), \lstinline|registerWindowControls| roteia cada IPC
de janela para \emph{a janela que o enviou}, via
\lstinline|BrowserWindow.fromWebContents(event.sender)|, e não para o singleton
\lstinline|state.mainWindow|. Os canais são: \lstinline|window:minimize|,
\lstinline|window:maximize-toggle|, \lstinline|window:close| e
\lstinline|window:get-state| (mais zoom: \lstinline|zoom-in|,
\lstinline|zoom-out|, \lstinline|zoom-reset|, e \lstinline|open-design-lab|).

\section{Ciclo de vida da aplicação: \file{main/lifecycle.js}}
\label{sec:05-lifecycle}

O módulo \file{main/lifecycle.js} concentra o ciclo de vida do app via a função
\lstinline|register()|, registrada por \file{main.js} no boot.

\subsection{Single-instance lock e múltiplas janelas}

No início, \lstinline|register()| detecta um \path{.spf} na linha de comando
(\lstinline|process.argv.find(arg => arg.endsWith('.spf'))|) e tenta adquirir o
\emph{single-instance lock} com \lstinline|app.requestSingleInstanceLock()|. A
variável de ambiente \lstinline|SAPHO_SKIP_SINGLE_INSTANCE=1| ignora o lock (uso
exclusivo dos testes). Se o lock não é obtido, a função loga o motivo em
\emph{ambos} o arquivo de log e o \lstinline|stderr| e chama
\lstinline|app.quit()|, retornando \lstinline|false| --- o que faz \file{main.js}
abortar o registro dos handlers num processo que está morrendo.

Quando uma segunda instância é lançada, o evento \lstinline|second-instance|
\textbf{cria uma nova janela} no mesmo processo (em vez de só focar a existente),
permitindo vários projetos lado a lado sem duplicar processos Electron nem correr
risco de corrida sobre a pasta compartilhada \path{components/Temp}. Se essa
segunda invocação trouxe um \path{.spf}, ele é enviado à nova janela.

\subsection{Encerramento: limpeza em duas fases}
\label{sec:05-before-quit}

O \emph{handler} \lstinline|before-quit| é a limpeza \emph{autoritativa} e roda em
duas fases serializadas, para evitar \lstinline|EBUSY| no \lstinline|rmdir| de
\path{Temp}:

\begin{description}[leftmargin=*]
  \item[Fase 1 --- soltar handles.] Fecha todos os \emph{watchers} de arquivo e
    de diretório (chokidar, que no Windows usa \lstinline|ReadDirectoryChangesW|)
    e chama \lstinline|stopAllToolchain()| (\autoref{sec:05-process-registry}),
    que mata todos os processos filhos da toolchain. Também encerra o servidor
    \tool{MCP} da Aurora (a ponte HTTP local). Tudo isso corre em paralelo, com
    um teto de \textbf{5\,s} via \lstinline|Promise.race| para não travar o quit.
  \item[Fase 2 --- apagar \path{Temp}.] Só depois de a fase 1 liberar os handles,
    remove recursivamente \path{components/Temp} (\lstinline|maxRetries:3|) e a
    recria vazia.
\end{description}

Há ainda \lstinline|window-all-closed| (sai do app fora do macOS) e
\lstinline|activate| (recria a janela principal se não houver nenhuma, padrão
macOS).

\section{Caminhos e estado do main}
\label{sec:05-paths-state}

\subsection{\file{main/paths.js}}

Centraliza os caminhos estáticos. \lstinline|appRoot| é
\lstinline|app.getAppPath()| (a raiz do projeto em dev, ou
\path{resources/app.asar} empacotado). O \lstinline|componentsPath| --- onde mora
o toolchain bundle (\autoref{cap:09-toolchain-bundle}) --- depende do modo: em dev
é \path{<appRoot>/components}; empacotado, fica ao lado do executável
(\lstinline|path.dirname(app.getPath('exe'))/components|), pois os componentes são
distribuídos \emph{fora} do asar. O sinalizador \lstinline|isDev| deriva de
\lstinline|process.env.NODE_ENV === 'development'|.

\subsection{\file{main/state.js}}
\label{sec:05-state}

Estado mutável compartilhado do \emph{main}. Todos os módulos importam a
\emph{mesma} referência de objeto, de modo que escritas feitas por um handler
ficam visíveis a todos. Substitui as antigas variáveis \lstinline|let| de escopo
de arquivo. Os campos incluem:

\begin{itemize}[leftmargin=*]
  \item \textbf{Janelas:} \lstinline|mainWindow|, \lstinline|splashWindow|,
    \lstinline|updateWindow|, \lstinline|prismWindow|.
  \item \textbf{Updater:} \lstinline|isQuitting|,
    \lstinline|downloadInProgress|, \lstinline|updateAvailable|,
    \lstinline|updateInfo| etc.
  \item \textbf{Projeto:} \lstinline|currentOpenProjectPath| (a fonte única da
    verdade sobre qual projeto está aberto, lida por \file{git.js} e
    \file{search.js}) e \lstinline|fileToOpen|.
  \item \textbf{Processos de simulação:} \lstinline|currentVvpProcess|,
    \lstinline|vvpProcessPid|, \lstinline|currentGtkwaveProcesses| e
    \lstinline|childProcesses| --- o \emph{Set} central de todo filho da
    toolchain, \emph{tree-killed} ao fechar a janela/sair.
  \item \textbf{Watchers:} \lstinline|activeWatchers|,
    \lstinline|activeDirectoryWatchers| e seus caches de \emph{stats}.
\end{itemize}

\section{Rastreio e \emph{tree-kill} de processos filhos}
\label{sec:05-process-registry}

O módulo \file{main/process_registry.js} é o registro central de todo processo
filho da toolchain que a \tool{AURORA} dispara (compiladores \cmm/ASM,
\tool{iverilog}, \tool{vvp}, \tool{Verilator} com seus \tool{g++}/\tool{make}/
\tool{ccache}, \tool{yosys} para o \tool{PRISM}, \tool{gtkwave}, \tool{cocotb}/
Python) mais a rotina única que força a parada de todos.

\subsection{Registro e ponto de spawn único}

\begin{description}[leftmargin=*]
  \item[\lstinline|trackChild(child)|] adiciona o filho a
    \lstinline|state.childProcesses| e o remove automaticamente em
    \lstinline|exit|/\lstinline|close|/\lstinline|error| --- o \emph{Set} só
    contém processos vivos. Marca também \lstinline|toolchainEverRan = true|.
  \item[\lstinline|spawnTracked(cmd, args, opts)|] é o ponto de \emph{spawn}
    obrigatório da toolchain: equivale a \lstinline|trackChild(spawn(...))|, com a
    mesma assinatura e o mesmo retorno de \lstinline|child_process.spawn|. Spawnar
    por aqui torna automático o registro para \emph{tree-kill}.
\end{description}

Os CLIs de agente de IA (\file{claude_code}, \file{codex_cli})
\emph{deliberadamente} não passam por aqui --- eles gerenciam as próprias árvores
de processo e são derrubados por seus \lstinline|killAll()| em
\lstinline|stopAllToolchain()|.

\subsection{\lstinline|stopAllToolchain()|}

É \emph{best-effort} e idempotente (memoizada numa única \lstinline|stopPromise|),
chamável tanto no \lstinline|close| da janela quanto no \lstinline|before-quit|.
Cobre, em ordem:

\begin{enumerate}[leftmargin=*]
  \item todo filho rastreado, \emph{tree-killed} por PID
    (\lstinline|taskkill /F /T /PID|);
  \item o processo de simulação \emph{legacy} de slot único, caso não estivesse
    rastreado;
  \item \textbf{apenas se algum filho realmente rodou nesta sessão}
    (\lstinline|toolchainEverRan|): varreduras de \emph{backstop} por nome
    (\lstinline|vvp.exe|, \lstinline|gtkwave.exe|) e por prefixo de caminho
    (os \lstinline|V<top>.exe| do Verilator sob \path{components/Temp});
  \item os CLIs de IA (Claude Code / Codex) via \lstinline|killAll()|;
  \item as gerações de chat de IA em curso, via \lstinline|abortAll()| (aborta os
    \emph{streams} HTTP).
\end{enumerate}

A guarda \lstinline|toolchainEverRan| é o que faz uma sessão só de edição fechar
instantaneamente: as varreduras externas (que custam um \lstinline|taskkill| frio
e um \lstinline|Get-CimInstance Win32_Process| via PowerShell) só rodam se houve
provadamente o que matar. As primitivas de \emph{kill} (\lstinline|killProcessSilently|,
\lstinline|killProcessesByName|, \lstinline|killProcessesByPathPrefix|) vivem em
\file{main/utils.js}.

\section{Higiene de temporários e log}
\label{sec:05-tempgc-log}

\subsection{\file{main/temp_gc.js}}

Higiene universal de temporários, \emph{best-effort} e não-bloqueante, rodada uma
vez no startup (enfileirada para fora do caminho crítico via
\lstinline|setImmediate|). \lstinline|runStartupGC()| (1) poda os arquivos
\path{aurora-mcp-<pid>.json} obsoletos que o Claude Code escreve no diretório
temporário do SO e que nunca eram limpos, e (2) limpa as imagens de anexo da IA
(\lstinline|cleanupTempImages(0)|), que são \emph{one-shot} por turno.

\subsection{\file{main/logger.js}}

Configuração central do
\href{https://github.com/megahertz/electron-log}{\tool{electron-log}}, que é um
singleton: como toda chamada \lstinline|require('electron-log')| no código retorna
a mesma instância, basta configurar transportes/níveis/formato uma vez aqui e todo
o \emph{main} herda. O transporte de arquivo fica no nível \lstinline|info|, com
rotação a $\sim$5\,MB; o transporte de console é \lstinline|debug| em dev e
\lstinline|warn| em produção. Os caminhos de log por SO (padrão do
\tool{electron-log}) são:

\begin{tabularx}{\linewidth}{Y Z}
  \toprule
  \textbf{Sistema} & \textbf{Caminho de \path{main.log}} \\
  \midrule
  Windows & \path{%APPDATA%/Aurora-IDE/logs/main.log} \\
  macOS   & \path{~/Library/Logs/Aurora-IDE/main.log} \\
  Linux   & \path{~/.config/Aurora-IDE/logs/main.log} \\
  \bottomrule
\end{tabularx}

\section{Carregando o renderer empacotado pelo Vite}
\label{sec:05-render-loader}

\subsection{\file{main/render_loader.js}}

Toda janela de primeira parte (principal, splash, atualização, \tool{PRISM})
carrega sua página por \lstinline|loadPage(win, relPath)|, centralizando a
seleção dev/prod num só lugar:

\begin{itemize}[leftmargin=*]
  \item \textbf{Dev (Vite):} quando \lstinline|AURORA_RENDERER_URL| está definida
    \emph{e} o app não está empacotado, a página é carregada do servidor de dev do
    Vite (com HMR), anexando o caminho relativo à origem do servidor.
  \item \textbf{Produção:} carrega a página empacotada de \path{dist/} via
    \lstinline|file://|.
\end{itemize}

Não há \emph{fallback} para a fonte crua: após a migração para \tool{Lit}, a
\path{index.html} crua carrega \lstinline|import|s de \lstinline|lit| que só
resolvem pelo \emph{bundler}, então um \path{dist/} ausente é um erro de build
--- e é exibido explicitamente como uma página de erro, não como uma UI
silenciosamente quebrada.

\subsection{\file{scripts/launch-electron.js}}

O \lstinline|npm start| não chama \lstinline|electron .| diretamente, e sim
\file{scripts/launch-electron.js}. O motivo: alguns shells pais (o terminal
integrado do VS Code, o Claude Code) exportam
\lstinline|ELECTRON_RUN_AS_NODE=1|; herdada por \lstinline|electron .|, ela faria
o Electron bootar como Node puro (\lstinline|app| viraria \lstinline|undefined| e
o app quebraria em \lstinline|app.getAppPath()| dentro de \file{main/paths.js}). O
script \emph{deleta} essa variável do ambiente e então faz
\lstinline|spawn(electron, ['.', ...argv])|.

\section{Topologia de IPC}
\label{sec:05-ipc}

A fronteira IPC é o coração da arquitetura: é por ela que o renderer alcança o
\emph{main}. No total existem \textbf{157} registros de \lstinline|ipcMain| no
processo principal (contagem por \lstinline|grep| de
\lstinline|ipcMain.handle/on/once| em \file{main/} e \file{main.js}). A maioria
usa \lstinline|ipcMain.handle| (\emph{request/response} via
\lstinline|invoke|/\lstinline|Promise|); alguns usam \lstinline|ipcMain.on|
(sinais unidirecionais, como controles de janela e a conclusão de \emph{tool
call}).

\subsection{Os módulos \file{main/ipc/*}}

\begin{longtable}{@{}p{0.18\linewidth} r p{0.58\linewidth}@{}}
  \toprule
  \textbf{Módulo} & \textbf{Handlers} & \textbf{Papel} \\
  \midrule
  \endhead
  \file{ipc/files.js}    & 27 &
    Sistema de arquivos, diálogos e \emph{watchers} (chokidar):
    \lstinline|read-file|, \lstinline|write-file|, \lstinline|copy-file|,
    \lstinline|mkdir|, \lstinline|file:delete|, \lstinline|watch-file|,
    \lstinline|watch-directory|, \lstinline|dialog:showOpen|,
    \lstinline|show-save-dialog|, \lstinline|create-backup| etc. \\
  \file{ipc/git.js}      & 28 &
    Controle de versão via \href{https://github.com/steveukx/git-js}{simple-git}
    sobre o diretório do projeto aberto: \lstinline|git:status|,
    \lstinline|git:diff|, \lstinline|git:commit|, \lstinline|git:push|,
    \lstinline|git:pull|, \lstinline|git:clone|, \lstinline|git:branches|,
    \lstinline|git:stash| etc. (\autoref{cap:11-git-github}). \\
  \file{ipc/ai.js}       & 22 &
    Aurora Intelligence: gestão de chaves (\lstinline|ai:set-key|,
    \lstinline|ai:get-key-status| --- a chave \emph{nunca} sai do main),
    \lstinline|ai:chat-start|/\lstinline|ai:chat-abort|, \emph{bridges} de
    Claude Code/Codex, manifesto de \emph{tools}, histórico de conversas
    (\autoref{cap:12-aurora-intelligence}). \\
  \file{ipc/project.js}  & 11 &
    Lifecycle de projeto (\lstinline|project:open|, \lstinline|project:close|,
    \lstinline|project:createStructure|) e CRUD de processadores
    (\lstinline|create-processor-project|, \lstinline|delete-processor|,
    \lstinline|rename-processor|); reescreve o \path{.spf} nesses eventos. \\
  \file{ipc/compile.js}  & 7 &
    Execução de compilação/simulação \emph{legacy} e abertura de visualizadores
    de waveform: \lstinline|launch-gtkwave-only|, \lstinline|launch-surfer|,
    \lstinline|cancel-vvp-process|, \lstinline|kill-current-spec-process|,
    \lstinline|check-process-running|, \lstinline|decode-complex|. \\
  \file{ipc/prism.js}    & 6 &
    \tool{PRISM}: dono da janela do visualizador de RTL e de todos os passos de
    síntese (Yosys $\to$ netlistsvg): \lstinline|prism-compile-with-paths|,
    \lstinline|prism-recompile|, \lstinline|generate-svg-from-module|,
    \lstinline|prism:build-digitaljs| (\autoref{cap:08-prism-rtl}). \\
  \file{ipc/github_auth.js} & 7 &
    Conexão de conta GitHub (PAT validado + \emph{Device Flow} OAuth), com o
    token cifrado por \lstinline|safeStorage| (DPAPI no Windows):
    \lstinline|github:connect|, \lstinline|github:status|,
    \lstinline|github:list-repos|, \lstinline|github:oauth-login| etc. \\
  \file{ipc/system.js}   & 6 &
    Utilidades de sistema: \lstinline|get-components-path|,
    \lstinline|toolchain:python-status|, \lstinline|join-path|,
    \lstinline|path-dirname|, \lstinline|app:reload| e
    \lstinline|renderer:error| (a \emph{error boundary} do renderer encaminha
    erros não tratados para o log persistente). \\
  \file{ipc/search.js}   & 2 &
    \emph{Find in Files} no projeto: \lstinline|search:in-project|, uma varredura
    recursiva síncrona limitada por profundidade/tamanho/matches, ancorada em
    \lstinline|state.currentOpenProjectPath|. \\
  \bottomrule
  \caption{Os nove módulos sob \file{main/ipc/} (somam 116 handlers).}
\end{longtable}

\subsection{Outros registradores de IPC}

Além de \file{main/ipc/*}, registram canais: \file{main/compile/executor.js}
(4 --- \lstinline|exec-spec|, \lstinline|exec-spec-streamed|,
\lstinline|exec-spec-allowed-binaries|, \lstinline|exec-spec-protected-flags|),
\file{main/lsp/verible_lsp.js} (9, \lstinline|lsp:*|),
\file{main/lsp/slang_lsp.js} (6, \lstinline|slang:*|),
\file{main/format/clang_format.js} (2, \lstinline|format:*|),
\file{main/treesitter/grammars.js} (2, \lstinline|treesitter:*| --- serve bytes
WASM ao realce semântico), \file{main/updater.js} (8, \lstinline|update*|) e
\file{main/windows.js} (10 --- os oito controles de janela/zoom mais os dois
sinais da splash \lstinline|splash:filled| e \lstinline|app:renderer-ready|,
registrados via \lstinline|ipcMain.once| dentro de \lstinline|createSplashScreen|).
Os LSP, \emph{clang-format} e \emph{tree-sitter} são todos
\emph{best-effort}: se o binário/grammar não estiver instalado, o lado \emph{main}
vira \emph{no-op}.

\section{Os preloads como ponte segura}
\label{sec:05-preloads}

Cada janela tem um preload que executa num contexto isolado e usa
\lstinline|contextBridge.exposeInMainWorld| para expor \emph{apenas} as funções
IPC que aquele renderer realmente consome --- uma allowlist de canais. Nenhum
preload expõe um \lstinline|send|/\lstinline|invoke| genérico, então um renderer
não pode alcançar um canal arbitrário.

\subsection{\file{js/app/preload.js} (janela principal)}

O preload principal expõe \emph{oito} \emph{namespaces} no
\lstinline|window| via \lstinline|contextBridge|:

\begin{tabularx}{\linewidth}{Y Z}
  \toprule
  \textbf{Objeto global} & \textbf{Conteúdo} \\
  \midrule
  \lstinline|electronAPI| & Arquivos, \emph{watchers}, projeto, compilação/
    simulação (\lstinline|execSpec|, \lstinline|execSpecStreamed|),
    \tool{PRISM}, diálogos, UI/janela, terminal, updater, busca,
    \lstinline|getPathForFile| e \lstinline|logRendererError|. \\
  \lstinline|terminalAPI| & Terminal integrado (canais do terminal expostos
    também como \emph{namespace} próprio, além do \emph{spread} em
    \lstinline|electronAPI|). \\
  \lstinline|aiAPI|       & Aurora Intelligence (provedores, chaves, chat,
    \emph{tools}, conversas). \\
  \lstinline|gitAPI|      & Git + integração GitHub. \\
  \lstinline|lspAPI|      & LSP Verilog (Verible). \\
  \lstinline|slangAPI|    & LSP semântico SystemVerilog (Slang). \\
  \lstinline|clangFormatAPI| & Formatação C/C++/\cmm. \\
  \lstinline|treeSitterAPI| & Bytes WASM de \emph{grammars} tree-sitter. \\
  \bottomrule
\end{tabularx}

Cada método é um \emph{wrapper} fino sobre \lstinline|ipcRenderer.invoke| ou
\lstinline|ipcRenderer.send|, mantendo o mapeamento 1:1 com os
\lstinline|ipcMain.handle/on| do \emph{main}. As chaves de API em texto puro
trafegam \emph{só} no sentido renderer $\to$ main: não existe canal
\lstinline|getKey|, pois ler bytes decifrados de volta para o renderer anularia o
propósito de manter o \emph{keystore} no \emph{main}. \lstinline|getPathForFile|
usa \lstinline|webUtils.getPathForFile| para resolver o caminho de um \lstinline|File|
em \emph{drag\&drop} (necessário no Electron 32+, que removeu \lstinline|file.path|).

\subsection{Preloads auxiliares}

\begin{description}[leftmargin=*]
  \item[\file{preload_splash.js}] expõe \lstinline|splashAPI| com apenas
    \lstinline|onProgress|, \lstinline|notifyFilled| e \lstinline|getAppVersion|
    --- a splash não precisa de mais nada.
  \item[\file{preload_update.js}] expõe \lstinline|updateAPI| espelhando 1:1 os
    canais de \file{main/updater.js}: \lstinline|onState|, \lstinline|onProgress|,
    \lstinline|onError|, \lstinline|onShake| (main $\to$ renderer) e
    \lstinline|download|, \lstinline|install|, \lstinline|dismiss|
    (renderer $\to$ main).
  \item[\file{preload_prism.js}] expõe um \lstinline|electronAPI| reduzido com
    cada canal que a janela do \tool{PRISM} usa (\lstinline|readFile|,
    \lstinline|prismCompileWithPaths|, \lstinline|buildDigitalJS|,
    \lstinline|openSourceFile| etc.), e comenta explicitamente que \emph{não} há
    \lstinline|send|/\lstinline|removeAllListeners| genéricos, para preservar a
    allowlist de canais.
\end{description}

\section{Diagrama da topologia}
\label{sec:05-diagrama}

O diagrama abaixo resume processos, janelas e a direção do IPC.

\begin{lstlisting}
+===========================================================================+
|                       PROCESSO MAIN (Node.js)                             |
|                                                                           |
|  main.js  -> lifecycle.register() (single-instance lock, before-quit)     |
|          -> <modulo>.register() para cada IPC abaixo                       |
|          -> windows.createSplashScreen() em app.whenReady                  |
|                                                                           |
|  Estado/infra:  state.js   paths.js   logger.js   temp_gc.js              |
|                 process_registry.js (trackChild / spawnTracked /          |
|                                      stopAllToolchain -> taskkill /T)      |
|                                                                           |
|  IPC (157 ipcMain):                                                        |
|    files  project  compile  git  github_auth  prism  ai  search  system   |
|    + executor(exec-spec)  lsp(verible/slang)  format  treesitter  updater |
|                                                                           |
|    spawn(shell:false) -> YANC / iverilog / vvp / verilator / yosys /      |
|                          gtkwave / cocotb   (todos rastreados)            |
+=================================|=========================================+
            contextBridge (preload)  |  ipcRenderer.invoke / .send  /  eventos
   +--------------+--------------+----+------------+-------------------+
   |              |              |                 |                   |
+--v---------+ +--v--------+ +---v---------+ +-----v--------+ +--------v------+
| RENDERER   | | RENDERER  | | RENDERER    | | RENDERER     | | RENDERER      |
| principal  | | splash    | | update      | | PRISM        | | Design Lab    |
| index.html | | splash.   | | update-     | | prism.html   | | design-lab.   |
|            | | html      | | notif.html  | |              | | html          |
| preload.js | | preload_  | | preload_    | | preload_     | | (sem preload) |
| (electron/ | | splash.js | | update.js   | | prism.js     | |               |
|  ai/git/   | | splashAPI | | updateAPI   | | electronAPI  | |               |
|  lsp/...)  | |           | |             | | (reduzido)   | |               |
+------------+ +-----------+ +-------------+ +--------------+ +---------------+
  ctxIso=on     ctxIso=on     ctxIso=on       ctxIso=on        ctxIso=on
  sandbox=on    sandbox=on    sandbox=on      sandbox=on       sandbox=on
  nodeInt=off   nodeInt=off   nodeInt=off     nodeInt=off      nodeInt=off
\end{lstlisting}

\begin{nota}
\textbf{Invariante de segurança.} Em todas as cinco janelas vale
\lstinline|contextIsolation:true|, \lstinline|sandbox:true| e
\lstinline|nodeIntegration:false|. O renderer nunca toca em Node, \lstinline|fs|
ou \lstinline|spawn| diretamente: tudo cruza a fronteira IPC por um canal
enumerado do preload, onde o \emph{main} valida a entrada
(\lstinline|safePath|, \emph{allowlist} de binários, \emph{protected flags}) antes
de executar a operação privilegiada.
\end{nota}
